% ####################################################################
% Basic configuration and packages
% ####################################################################
% Package for discovering wrong and outdated usage of LaTeX.
% More information to be found in l2tabu English version.
\RequirePackage[l2tabu, orthodox]{nag}
% Class of LaTeX document: {size of paper, size of font}[document class]
\documentclass[a4paper,11pt]{article}



% ############################
% Packages not tied to particular normal language
% ############################
% This package should improved spaces in the text
\usepackage{microtype}
% Add few important symbols, like text Celcius degree
\usepackage{textcomp}



% ##################
% Polonization of LaTeX document
% ##################
% Basic polonization of the text
\usepackage[MeX]{polski}
% Switching on UTF-8 encoding
\usepackage[utf8]{inputenc}
% Adding font Latin Modern
\usepackage{lmodern}
% Package is need for fonts Latin Modern
\usepackage[T1]{fontenc}



% ##################
% Setting margins
% ##################
\usepackage[a4paper, total={14cm, 25cm}]{geometry}



% ##################
% Setting vertical spaces in the text
% ##################
% Setting space between lines
\renewcommand{\baselinestretch}{1.1}

% Setting space between lines in tables
\renewcommand{\arraystretch}{1.4}



% ##################
% Packages for scientific papers
% ##################
% Switching off \lll symbol, that I guess is representing letter "Ł"
% It collide with `amsmath' package's command with the same name
\let\lll\undefined
% Basic package from American Mathematical Society (AMS)
\usepackage[intlimits]{amsmath}
% Equations are numbered separately in every section
\numberwithin{equation}{section}

% Other very useful packages from AMS
\usepackage{amsfonts}
\usepackage{amssymb}
\usepackage{amscd}
\usepackage{amsthm}

% Package with better looking calligraphy fonts
\usepackage{calrsfs}

% Package with better looking greek letters
% Example of use: pi -> \uppi
\usepackage{upgreek}
% Improving look of lambda letter
\let\oldlambda\Lambda
\renewcommand{\lambda}{\uplambda}




% ##################
% BibLaTeX
% ##################
% Package biblatex, with biber as its backend, allow us to handle
% bibliography entries that use Unicode symbols outside ASCII
\usepackage[
language=polish,
backend=biber,
style=alphabetic,
url=false,
eprint=true,
]{biblatex}

\addbibresource{Informatyka-teoretyczna-Bibliography.bib}





% ##################
% Defining new environments (?)
% ##################
% Defining enviroment "Wniosek"
\newtheorem{corollary}{Wniosek}
\newtheorem{definition}{Definicja}
\newtheorem{theorem}{Twierdzenie}





% ##################
% Local packages
% You need to put them in the directory Local-packages
% ##################
% Local configuration of this particular file
\usepackage{./Local-packages/local-settings}

% Package containing various command useful for working with a text
\usepackage{./Local-packages/general-commands}
% Package containing commands and other code useful for working with
% mathematical text
\usepackage{./Local-packages/math-commands}





% ######################################
% Local macros
% ######################################
\newcommand{\rel}{\text{rel}}

\newcommand{\twoAngels}{\langle \hspace{1em} \rangle}

\newcommand{\doKeyword}{\text{\textbf{do}}}
\newcommand{\untilKeyword}{\text{\textbf{until}}}
\newcommand{\whileKeyword}{\text{\textbf{while}}}




% ######################################
% Package "hyperref"
% They advised to put it on the end of preambule
% ######################################
% It allows you to use hyperlinks in the text
\usepackage{hyperref}










% ####################################################################
% Title and author of the text
\title{Informatyka teoretyczna \\
  {\Large Błędy i~uwagi}}

\author{Kamil Ziemian \\
  \email}


% \date{}
% ####################################################################










% ####################################################################
% Beginning of the document
\begin{document}
% ####################################################################





% ######################################
% Title of the text
\maketitle
% ######################################





% ######################################
\section{J.M. Brady \textit{Informatyka teoretyczna w~ujęciu
    programistycznym},
  \parencite{Brady-Informatyka-teoretyczna-w-ujeciu-ETC-Pub-1983}}

\label{sec:Brady-Informatyka-teoretyczna-ETC}
% ######################################


% ##################
\CenterBoldFont{Uwagi do~konkretnych stron}


\noindent
\StrWierszDol{290}{15} Ten wiersz, i~inne podobne do niego, wyglądałby
lepiej, gdyby zamiast „//zbiór liczb parzystych” był zapisany jako
„//~zbiór liczb parzystych”.

\VerSpaceFour





\noindent
\StrWierszeDol{292}{3--4} Czytamy tutaj „Nazwa \textit{zbiór potęgowy}
wywodzi~się nie tyle z~potęgi takiej operacji, ile ze spostrzeżenia, jakie
poczynimy w~poniższym ćwiczeniu.”. Angielska nazwa tego zbioru to
\textit{power set}, więc w~oryginalne tekst ten brzmiał prawdopodobnie jakoś
tak: „The name \textit{power set} comes not from the~power of this
operation\ldots”. Cały żart w języku angielskim brzmi lepiej, niż po polsku.

\VerSpaceFour





\noindent
\StrWierszDol{293}{1} Wzór zawarty w~tej linii
\begin{equation}
  \label{eq:Brady-Infromatyka-ETC-01}
  f | U ( u ) = f( u ) \in T,
\end{equation}
byłby łatwiejszy w~zrozumieniu, gdyby go zapisać jako
\begin{equation}
  \label{eq:Brady-Infromatyka-ETC-02}
  \big( f | U \big)( u ) = f( u ) \in T,
\end{equation}

\VerSpaceFour





\noindent
\StrWierszDol{295}{6} W~miejscach takich jak to, widać, że~w~tym wydaniu
funkcja~$g$ obliczana na wartości~$t$ ma zapis graficzny „$g \, ( t )$”.
Ładniej by wyglądała, gdyby była zapisywana jako~„$g( t )$”.

\VerSpaceFour





\noindent
\StrWierszGora{296}{10} Do zapisu relacji wprowadzony jest tu symbol
$- \times$. Można zgadywać, że~tym czego naprawdę pragnął autor, jest strzałka,
która kończy~się znakiem~$\times$. Można spróbować dodać taką strzałkę
do~\LaTeX a, jeśli nie jest dostępna w~którejś z~istniejących paczek.

\VerSpaceFour





\noindent
\StrWierszGora{296}{15} Zostaje tu wprowadzony symbol $rel_{ U }$,
oznaczający relację zdefiniowaną przez zbiór $U \subseteq S \times T$. Wyglądałby on
jednak lepiej, gdyby był on zapisywany jako $\rel_{ U }$.

\VerSpaceFour





\noindent
\StrWierszGora{296}{17} W~tej linii spotykamy wzór $\rel_{ U }( s ) = t$,
którego sens jest zbyt dwuznaczny. Może on bowiem sugerować, że~$t$ jest
\textit{jedynym} elementem, który jest w~relacji $\rel_{ U }$
z~elementem~$s$. Jednak to nie jest to co ten wzór ma oznaczać, gdyż mówi
on tylko, że~$s$ i~$t$ są w~relacji, co wynika z~tego,
że~$\rel_{ U }( s ) = t$ wtedy i~tylko wtedy, gdy $( s, t ) \in U$. Z~tego
powodu symbol ten lepiej by było używać w~innej, rozpowszechnionej
w~literaturze notacji $s \, \rel_{ U } \, t$.

\VerSpaceFour





\noindent
\StrWierszGora{297}{18} Funkcje o~nazwach takich jak wprowadzona tu funkcja
„$nast$”, lepiej byłoby chyba zapisywać jako „$\text{nast}$”.

\VerSpaceFour





\noindent
\StrWierszGora{301}{2} W~podanym tu zapisie liczby rzeczywistej
$n \cdot d_{ 1 } d_{ 2 } d_{ 3 }\ldots$, kropka „$\cdot$” oznacza, że~po jej lewej stronie
znajduje~się symbol liczby całkowitej, a~po jej prawej stronie cyfry
jej części ułamkowej. Jak tekst wyjaśnia, używana jest tutaj reprezentacja
dziesiętna, czyli $d_{ i } \in \{ 0, 1, \ldots, 9 \}$.

Autor przeczył chyba przy tym fakt, że~skoro przyjął $n \in \Nbb$, to jego
definicji obejmuje tylko nieujemne liczby rzeczywiste. Jest to jednak tylko
drobne potknięcie, bez większego znaczenia.





\noindent
\StrWierszGora{303}{5} Warto doprecyzować, że~w~wyrażenie takie jak
$1/2 * n * ( n + 1 )$ należy rozumieć jako $( 1 / 2 ) * n * ( n + 1 )$.

\VerSpaceFour





\noindent
\StrWierszGora{307}{5} W~tej książce przyjęto konwencję, że~kolejne elementy
ciągu zapisujemy jako $e1, e2, e3, \ldots$, nie zaś jako
$e_{ 1 }, e_{ 2 }, e_{ 3 }, \ldots${ } Od konwencji tej postanowiono odejść,
zapisują wyraz o~numerze $n + 1$ jako $e_{ n + 1 }$, jednak dla spójności
zapisu proponujemy, by trzymać~się uprzedniej i~zapisywać ten wyraz jako
$e( n + 1 )$.

\VerSpaceFour





\noindent
\StrWierszGora{309}{19} Instrukcję iteracyjną
\begin{equation}
  \label{eq:Brady-Infromatyka-ETC-01}
  \untilKeyword \;\, ident > v2 \;\, \doKeyword \;\, instrukcja \;\,
  \twoAngels \;\, ident := ident + 1,
\end{equation}
można prościej zapisać jako
\begin{equation}
  \label{eq:Brady-Infromatyka-ETC-01}
  \whileKeyword \;\, ident \leq v2 \;\, \doKeyword \;\, instrukcja \;\,
  \twoAngels \;\, ident := ident + 1,
\end{equation}

\VerSpaceFour





\noindent
\StrWierszGora{311}{12} Ta pozycja bibliograficzna kończy~się jako
„\textit{languages, maszyn}”, co wskazuje, że~została błędnie wpisana.

\VerSpaceFour






% ##################
\newpage

\CenterBoldFont{Błędy}


\begin{center}

  \begin{tabular}{|c|c|c|c|c|}
    \hline
    Strona & \multicolumn{2}{c|}{Wiersz} & Jest
    & Powinno być \\ \cline{2-3}
           & Od góry & Od dołu & & \\
    \hline
    % & & & & \\
    % & & & & \\
    % & & & & \\
    % & & & & \\
    290 & & 13 & liczb rzeczywistych & dodatnich liczb wymiernych \\
    291 & & \hphantom{0}1 & $\varepsilon$. & $\varepsilon$.” \\
    299 & \hphantom{0}9 & & dyskusjiz & dyskusji z \\
    303 & & 11 & $1 \ldots n$ & $1 + \ldots + n$ \\
    311 & 11 & & in~Programming & \textit{in~Programming} \\
    311 & & \hphantom{0}2 & theory \textit{Scientific}
           & theory, \textit{Scientific} \\
    313 & 16 & & \textit{Springer} & Springer \\
    314 & & \hphantom{0}7 & (1974):\textit{Mathematical}
           & (1974): \textit{Mathematical} \\
    315 & 17 & & \textit{Mass. Inst. Technol.}
           & Mass. Inst. Technol. \\
    316 & 15 & & \textit{languages} & \textit{Languages} \\
    316 & 17 & & 150.{ }{ } Scott & 150. Scott \\
    316 & & \hphantom{0}4 & Thue{ }{ } A. & Thue A. \\
    316 & & \hphantom{0}2 & Top\hspace{0.2em}or & Topor \\
    \hline
  \end{tabular}

\end{center}

\VerSpaceSix





\noindent
\StrWierszeDol{310}{6--8} \\
\Jest Na przykład poniższa funkcja daje w~wyniku albo pierwszą wartość
większą od $m$, wartości wybranej ze~struktury danych $V$ o~$k$ elementach,
albo też $0$, jeśli nie ma w~strukturze takiej wartości \\
\PowinnoByc Dla przykładu rozpatrzmy poniższą funkcję, która operuje
na strukturze zawierającej $k$~elementów. Funkcja ta zwraca albo pierwszą
wartość zawartą w~tej strukturze, która jest większa od~$m$, albo też~$0$,
jeśli struktura nie zawiera takiej wartości. \\


% ####################################################################










% ####################################################################
% ####################################################################
% Bibliography

\printbibliography





% ############################
% End of the document

\end{document}
